\documentclass[a4paper,12pt]{article}
\usepackage[margin=2cm]{geometry}

% Pacchetti utili di base
\usepackage[utf8]{inputenc}   % codifica dei caratteri
\usepackage[T1]{fontenc}      % font encoding
\usepackage[italian]{babel}   % lingua italiana
\usepackage{graphicx}         % immagini
\usepackage{hyperref}         % link cliccabili
\usepackage{amsmath}          % align center
\usepackage{listings}         %codice programmazione
\usepackage{xcolor}           % per colori opzionali del codice

% Informazioni sul documento
\title{Riassunto per lo scritto di T.W.}
\author{Nicolò Giuliani}
\date{}

\lstset{
  language=HTML,
  basicstyle=\ttfamily\small,
  keywordstyle=\color{blue},
  stringstyle=\color{red},
  commentstyle=\color{gray},
  breaklines=true
}
\begin{document}
\maketitle
\section{HTTP}
La richiesta via HTTP e'composta:
\begin{itemize}
    \item \textit{Method}: azione del server che deve eseguire (che il client richiede)
    \item \textit{URI}: risorsa locale del server
    \item \textit{Version}
    \item \textit{Header}: sono i metadati, divisi in \textbf{generali} (protocollo, timestamp, ecc), \textbf{entita'} (Content-Type, Content-Length, Content-Encoding, ...), \textbf{richiesta} (GET/POST)
    \item \textit{Body}: messaggio MIME (plain text / HTML / Base64 / ...)
\end{itemize}
Vediamo i metodi di richiesta:
\begin{itemize}
    \item GET: ottenere una risorsa (link o ricerca nel browser)
    \item HEAD: uguale a GET ma contiene solo header (senza body), serve solo per la validazione delle risorse
    \item POST: trasmettere informazioni al server rispetto alla risorsa del URI
    \item PUT: trasmettere informazioni al server creando o sostituendo la risorsa specificata (inverso del GET)
    \item DELETE
    \item PATCH: aggiornare parzialmente una risorsa
    \item OPTIONS: verificare opzioni, requisiti e servizi
\end{itemize}
La relativa risposta e'divisa in:
\begin{itemize}
    \item \textit{Status code} della richiesta
    \item Version
    \item Header
    \item Body: messaggio MIME
\end{itemize}
\section{HTTP/1.1}
Prima miglioria di speedup:
\begin{itemize}
    \item \textit{pipeline}: invio di piu'richieste senza aspettare risposta, ma le risposte arrivano in ordine delle richieste
\end{itemize}
\section{HTTP/2}
Migliorie delle performance tra cui
\begin{itemize}
    \item \textit{multiplexing} (parallelalizzazione delle richieste/risposte ora \textbf{compresse} (30\% di speedup))
    \item \textit{push anticipato} del server per il client
\end{itemize}
\section{HTTP/3}
Passaggio dal protoccollo viaggente su TCP/IP (HTML 1/2) ad QUIC (basato su UDP). \\
Dato che HTTP/x e'stateless per il server (non mantiene info delle sessioni precedenti), il clienti memorizza delle infomazioni delle scorse sessioni (\textbf{cookies}). Questi sono composti:
\begin{itemize}
    \item nome
    \item valori
    \item \textit{domain}: il dominio nel quale e'valido
    \item \textit{path}: il percorso del dominio nel quale e'valido
    \item \textit{expire}: quando non e'piu'valido
    \item \textit{secure}: se il client deve utilizzare \textbf{solo} HTTPS
    \item \textit{version}
\end{itemize}
Come si utilizzano?
\begin{itemize}
    \item \textit{cookies persistenti}: validita'lunga o non scadono, tipo login
    \item \textit{cookies sessione}: breve durata (cancellati alla chiusura del browser) per raggrupare operazioni in sessioni di lavoro
    \item \textit{cookies di terze parti}: esterne al dominio della pagina, tipo banner pubblicitari
\end{itemize}
\section{URI}
Sono la fusione degli URL (link) e URN (nome univoco come stringa, tipo isbn di un libro). \\
Gli URI sono progettati per:
\begin{itemize}
    \item Trascrivibili: hanno un set di caratteri limitato (cosi possono essere scritti cartelloni, carta, ecc)
    \item Identificazione: i protocolli non sono definiti dal URI
\end{itemize}
\begin{center}
    URI = schema : [//authority] path [?query][\#fragment]
\end{center}
Analizziamo:
\begin{itemize}
    \item Schema: protocollo TCP
    \item Authority: se presente indica gerarchia dei nomi ([userinfo @] host [: port])
    \item Path: percorso della risorsa
    \item Query: ulteriore specificazione della risorsa (generalmente per risposte dinamiche)
    \item Fragment: individua risorsa secondaria (un frammento) rispetto alla primaria
\end{itemize}
\subsection{Route}
Fatta una richesta di una risorsa il server restituisce in base ad:
\begin{itemize}
    \item Managed route: hardcoded ogni richiesta da una specifica risorsa (puo'essere anche dinamica)
    \begin{lstlisting}
    GET(http://www.example.com/name) --> Nicolo'Giuliani
    GET(http://www.example.com/css/style.css) --> raw del file css/style.css
    \end{lstlisting}
    \item File-system route: la radice del server corrisponde alla radice dell'URI
    \begin{lstlisting}
    www/
        index.html       --> http://www.example.com/
        alice/
            index.html   --> http://www.example.com/alice/
        css/
            style.css   --> http://www.example.com/css/style.css
    \end{lstlisting}
\end{itemize}
\subsection{Protocolli morti}
\begin{itemize}
    \item File: file://host/path [\#fragment] (dal browser)
    \item Data: data:[<media type>][;base64],<data> (non cerca ma CONTIENE la risorsa, puo'essere ANCHE codificata in Base64)
    \item FTP: ftp://[user[:password]@]host[:port]/path
\end{itemize}
\section{API REST}
Servizio basato su HTTP e URI come interfaccia per un'applicazione. \\
Si dice RESTFULL (o anche REST) quando e'composta:
\begin{itemize}
    \item Route (URI): collezioni o singoli elementi
    \item medoti HTTP (GET/POST/PUT/DELETE)
    \item input/output
    \item errori e relativi messaggi da restituire
\end{itemize}
La struttura moderna e'il \textbf{modello CRUD}:
\begin{itemize}
    \item Create
    \item Read
    \item Update
    \item Delete
\end{itemize}
\subsection{HTTP Caching}
Nelle applicazioni HTTP RESTFULL (API con molte richieste) possiamo avere una cache-server che velocizza la computazione del tempo di risposta delle richieste! NON riduce il traffico sulla rete, per quello c'e'altro. \\
Per essere sicuri che un oggetto in cache sia ancora valido:
\begin{itemize}
    \item \textit{Server-specified expiration}: scadenza della risorsa
    \item \textit{Heuristic expiration}: durata euristica del oggetto in base ad alcune info, tipo "Last-Modified time" 
\end{itemize}
\section{Base64}
Server per spedire binario via protocolli che vogliono testo (ASCII o UTF-8). \\
Metodo di codifica e decodifica (algoritmica banale) che permette di codificare a gruppi di 3byte in blocchi da 4 caratteri Base64.
\begin{center}
    3byte $\to$ 4 caratteri Base64 (4byte) \\
    4byte $\to$ 8 caratteri Base64 (8byte) \\
    5byte $\to$ 8 caratteri Base64 (8byte) \\
    6byte $\to$ 8 caratteri Base64 (8byte) 
\end{center}
\section{JSON Web Token}
Token leggero per autenticare client $\to$ server.
\begin{center}
    header.payload.signature
\end{center}
\begin{itemize}
    \item \textit{header}: tipologia (JWT) e algoritmo (SHA256)
    \item \textit{payload}: contiene le informazioni di scambio (claims)
    \item \textit{signature}: garantisce che il token non e'stato modificato, cosi:
    \begin{center}
        base64(header) + "." + base64(payload)
    \end{center}
\end{itemize}
\section{CORS}
Se in un sito abbiamo un codice maligno che "sniffa" i dati e li rigira ad un altro sito, abbiamo un \textbf{cross‑domain request}. \\
CORS e'un meccanismo HTTP che limita le richieste da solo determinati domini esterni! (Funziona per Ajax non per HTML classico).
\section{HTML}
Protoccollo client-server non solo per pagine Web ma anche per appliazioni con richiesta-risposta (API).
\subsection{Attributi globali}
\begin{itemize}
    \item \textit{id}: identificati univoco del oggetto
    \item \textit{class}: la classe del oggetto
    \item \textit{style}: stile grafico del oggetto
    \item \textit{title}: info extra per accessibilita'
\end{itemize}
\subsection{Lunghezze}
\begin{itemize}
    \item width = "150";
    \begin{itemize}
        \item valore assoluto di pixel
    \end{itemize}
    \item width = "100\%";
    \begin{itemize}
        \item valore proporzionale al contenitore
    \end{itemize}
    \item width = "150, 25\%, *, *";
    \begin{itemize}
        \item per le tabelle serve per distribuire le varie colonne come:
        \begin{itemize}
            \item 150px
            \item 25\% della pagina
            \item restante delle pagina / 2
            \item restante delle pagina / 2
        \end{itemize}
    \end{itemize}
\end{itemize}
\section{HTML5}
Base del web moderno, con media audio/video nativi, API complesse e nuovi elementi semantici.
\begin{itemize}
    \item <video> / <audio>: controlli nativi, piu'sorgenti e sottotitoli. Non servono piu'plugin esterni!
    \item <header> / <footer>: Intestazione, menu, logo, titolo, copyright, autori, ecc
    \item <article: contenuto autonomo (blog)
    \item <section>: contenuti raggrupati per argomento
    \item <nav>: menu, link pagine, breadcrumb
    \item <aside>: extra come sidebar, pubblicita, note
    \item <canvas>: grafica 2D via javascript
\end{itemize}
Codice piu'leggeibile e \textbf{SEO-friendly}.
\section{CSS}
Sintassi:
\begin{lstlisting}
    //per un elemento di HTML
    selettore { 
        [proprieta']: [statement];
    }
    //per una classe da noi definita
    .classe{

    }
    //per un oggetto UNICO
    #id{

    }
    //per la prima riga di p
    p:first-line{ }

    //per solo il testo di p fra "<<...>>"
    p{
        color: red;
    }
    q::before { content :"<<";}
    q::after { content:">>";}
\end{lstlisting}
HTML prevede l'uso di CSS in 3 modi:
\begin{itemize}
    \item presso il TAG di rifermento
    \item nel <style> dell'head
    \item via file esterno con <link>
\end{itemize}
E si assegna agli elementi con:
\begin{itemize}
    \item per tipo (nome elemento)
    \item per categoria (class)
    \item specifico elemento (id)
\end{itemize}
Gli attibuti vengono presi a \textbf{cascata} ovvero
\begin{center}
    \begin{lstlisting}
            p { font-family: Arial; font-size: 12 pt; }
            p { color: red; font-size: 11 pt; }
            p { margin-left: 15 pt; color: green;}
    \end{lstlisting} 
    viene trasformato in \\
    \begin{lstlisting}
            p {
                font-family: Arial;
                font-size: 11 pt;
                margin-left: 15 pt;
                color: green;
            }
    \end{lstlisting}
\end{center}
\subsection{Priorita'CSS}
\begin{itemize}
    \item Caso normale:
\begin{lstlisting}
    p { color: red; }
    p { color: green; }
\end{lstlisting}
avremo 
\begin{lstlisting}
    p{
        color: green;
    }
\end{lstlisting}
    \item Caso \textbf{!important}:
\begin{lstlisting}
    p { color: red; !important}
    p { color: green; }
\end{lstlisting}
avremo 
\begin{lstlisting}
    p{
        color: red;
    }
\end{lstlisting}
\item In CSS, quando più regole si applicano allo stesso elemento, vince quella con \textbf{specificità maggiore}. I selettori sono ordinati per importanza in base a questo principio: selettori più specifici coprono quelli più generali.
\begin{lstlisting}
<style>
    nav#main-nav li p { color:red; }
    #main-nav .active p { color:yellow; }
    body nav#main-nav li p { color:blue; }
    nav ul li.active p { color:green; }
</style>

<nav id="main-nav">
  <ul>
    <li class="active">
      <p style="color:black;">link</p>
    </li>
  </ul>
</nav>
\end{lstlisting}
\begin{table}[h!]
\centering
\begin{tabular}{|l|c|l|}
\hline
\textbf{Regola} & \textbf{Specificità} & \textbf{Commento} \\
\hline
nav\#main-nav li p & 0-1-0-3 & ID=1, classi=0, tag=3 \\
\#main-nav .active p & 0-1-1-1 & ID=1, classi=1, tag=1 \\
body nav\#main-nav li p & 0-1-0-4 & ID=1, classi=0, tag=4 \\
nav ul li.active p & 0-0-1-4 & ID=0, classi=1, tag=4 \\
\texttt{Inline} & \textbf{1-0-0-0} & Stile scritto direttamente nell'elemento, specificità massima \\
\hline
\end{tabular}
\vspace{1em}\\
$\to$ "color: black" 
\caption{Confronto di specificità dei selettori CSS incluso lo stile inline}
\end{table}
Riassumendo:
\begin{center}
    Inline > ID > class > TAG (elemento HTML)
\end{center}
\end{itemize}
\section{Javascript}
Durante l'esecuzione dello script, il browser è \textbf{bloccato} e non reagisce agli input dell'utente. \\
Come si esegue:
\begin{itemize}
    \item Sincrono:
    \begin{itemize}
        \item Appena lo script viene letto da <script> o da file
    \end{itemize}
    \item Asincrono:
    \begin{itemize}
        \item \textbf{event-oriented processing} (ad un evento, es. click bottone)
        \item aspettando una operazione di rete
        \item dopo un timeout il codice viene eseguito automaticamente
    \end{itemize}
\end{itemize}
Abbiamo diverse opzioni di output
\begin{itemize}
    \item console.log();
    \item alert();
    \item document.write(); (scrive nel documento direttamente)
    \item document.getElementById(id).innerHTML = string; (modificare il DOM visualizzato)
\end{itemize}
Tipi di dato:
\begin{itemize}
    \item bool
    \item numeri (int e float)
    \item string
    \item null
    \item undefined
    \item \textbf{object} (strutture e array)
    \begin{lstlisting}
        oggetto = {
            chiave: valore;
            "altrachiave": altrovalore;
        }
    \end{lstlisting}
    $\to$ tutte le chiavi sono trattate come stringhe!
\end{itemize}
\subsection{Accesso alle proprietà di un oggetto in JavaScript}
\begin{lstlisting}
const utente = {
  "profilo": {
    "info.personali": {
      "data-nascita": "1990-01-01"
    }
  }
};
\end{lstlisting}
Possiamo utilizzare la notazione \textit{punto} solo se non abbiamo caratteri come spazi, trattini o punti:
\begin{center}
    \begin{lstlisting}
        utente.profilo
    \end{lstlisting}
\end{center}
Mentre per le chiavi con simboli strani facciamo:
\begin{center}
    \begin{lstlisting}
        utente.profilo.["info.personali"];
    \end{lstlisting}
\end{center}
\section{Json}
JSON (JavaScript Object Notation) è un formato dati derivato dalla notazione usata da JS per gli oggetti.
\begin{lstlisting}
{
  "nome": [ "Giuseppe", "Andrea", "Federico" ],
  "cognome": "Rossi",
  "altezza": 180,
  "nascita": "1995-04-11T22:00:00.000Z",
  "indirizzo": {
    "via": {
      "strada": "Via Indipendenza",
      "numero": "15"
    },
    "citta": "Bologna",
    "nazione": "Italia"
  },
  "telefono": [
    { "tipo": "casa", "numero": "051  123456" },
    { "tipo": "cell", "numero": "335  987654" }
  ]
}
\end{lstlisting}
Oggigiorno ogni linguaggio di programmazione decente ha un modo per intefacciarsi con i json.
\section{Il Document Object Model (DOM)}
E'una API per i documetni HTML e XML. \\
Ogni componente del documento puo essere
\begin{itemize}
    \item letto
    \item modificato
    \item cancellato
\end{itemize}
dal DOM. \\
\vspace{5px}\\
Javascript implementa i metodi standard per accedere al DOM del documento.
\begin{lstlisting}
    var c = document.getElementById('c35');
    c.setAttribute('class','prova1');
    c.removeAttribute('align');
    var newP = document.createElement('p');
    var text = document.createTextNode('Ciao Mamma.');
    newP.appendChild(text);
    c.appendChild(newP);
\end{lstlisting}
\subsection{Selettori base HTML in DOM}
\begin{itemize}
    \item getElementById
    \item getElementsByName
    \item getElementsByTagName 
\end{itemize}
\section{Typescript}
E' un super-set di JavaScript, perciò ogni costrutto valido in JavaScript è valido in Typescript. \\
Vengono aggiunti costrutti per l’annotazione dei tipi e un sistema di type-inference che consente di conoscere (staticamente) il tipo di una variabile.
\section{Node.js}
Node.js è un ambiente di esecuzione Javascript progettato per costruire applicazioni server-side efficienti. \\
Invece di dedicare un thread ad ogni connessione ricevuta dalla rete, l'applicazione nodeJS utilizza un solo thread per le richieste (asincrone e non bloccanti) e tutte ricevono attenzione condivisa.
\begin{itemize}
    \item Mantenendo tutto in un unico stack, NodeJS può gestire decine o centinaia di migliaia di richieste concorrenti. (rispetto a poche migliaia)
\end{itemize}
\section{Express.js}
Express è un modello di server web per node.js:
\begin{itemize}
    \item open-source
    \item molto diffuso, ottima comunity e plugin
\end{itemize}
Express si dedica al routing, alla gestione delle sessioni, all'autenticazione degli utenti, e molti altri principi fondanti delle connessioni HTTP.
\end{document}